\section{Evidence Contract}

The sandbox compiles two disjoint objects.  A \emph{scenario} $s$ freezes model
and tokenizer revisions, token lengths, request count, arrival regime, target
hardware, and SLO.  A \emph{configuration} $x\in\mathcal X$ contains only
deployable controls: runtime, precision, batching and cache limits, TP/PP/DP,
placement, and power cap.  This prevents a search policy from reporting an
``improvement'' by silently reducing the requested work.

Every backend returns an append-only record
\begin{equation}
e=\bigl(h(s),h(x),\ell,k,\mathbf y,u,c,z\bigr),
\end{equation}
where $h(\cdot)$ is a SHA-256 identity, $\ell$ is fidelity, $k$ distinguishes
predicted and measured records, $\mathbf y$ holds energy and SLO metrics, $u$
is declared uncertainty, $c$ is measured GPU-hour cost, and $z$ records
software, topology, telemetry, and failure provenance.  L0 is a deterministic
CPU projection; L1 is a phase-controlled single-GPU measurement.  A
topology-aware CPU projection remains \emph{extrapolated} and retains
\texttt{validation\_required}.  Zero incremental GPU-hours is not zero cost:
CPU time, memory, and energy remain reportable overheads.

Predictions, intervals, and manifests are frozen before a holdout.  The
validator checks timestamps, hashes, split labels, scenario/configuration,
container identity, power readback, repeat count, and every raw artifact.  A
comparison is rejected if any check fails.  The present study exercises L1
calibration, L0 CPU prediction, and independent L1 validation; it makes no
topology-aware or target-scale accuracy claim.
